\section{Spektralanalyse}

In dieser Messreihe haben wir die unterschiedlichen Spektren des GEM-Detektors vermessen, wenn er verschiedenen Quellen ausgesetzt ist. Wir haben die Quellen $^{55}$Fe, $^{109}$Cd und die x-ray gun verwendet. \\ 
Zur Analyse der aufgenommenen Spektren haben wir zuerst an die Peaks der einzelnen Spektren Gaußkurven gefittet. Die aus diesem Fit erhaltenen Erwartungswerte und deren Unsicherheiten sind in Tabelle \ref{tab:Spektralanalyse} aufgeführt. Die Unsicherheiten, die wir dem Programm übergeben haben, setzen sich zusammen aus der quadratischen Summe des Poissonfehlers und des systematischen Fehlers aufgrund der Bingrößen: 


\begin{equation}
    \boxed{ \sigma_{fit} = \sqrt{N + \left(\frac{\text{Bingröße}}{2\sqrt3}\right)^2}}
    \label{eq:sigma_mu}
\end{equation}

N bezeichnet die counts im zugehörigen pulseheight-Intervall. Dabei wird ein Poission-Fehler auf die Counts angenommen. \\

Die erhaltenen Spektren und Gaußkurven sind in Abbildung \ref{fig:Spektren} zu sehen. 


\begin{figure}[h]
    \centering
    \includesvg[width=1\textwidth]{Bilder/Spektren.svg}
    \caption{Mit dem GEM-Detektor vermessene Spektren von Eisen, Cadmium und der x-ray gun. An jeden Peak wurde eine Gaußfunktion angepasst und der Peak bestimmt}
    \label{fig:Spektren}
    
\end{figure}


Die Peaks des $^{55}$Fe sind bekannt. Der Argon-Escape-Peak befindet sich bei \SI{2.95}{\keV} und der Peak des Eisens befindet sich bei \SI{5.88}{\keV} \cite{Vorlesungsfolien}. Vergleicht man diese Werte mit den im vermessenen Spektrum erhaltenen pulseheights ergeben sich die Kalibrierungsfaktoren und deren Unsicherheiten zu: \\
$K=\frac{\SI{5.88}{\keV}}{246.822} = 0.024 \pm \SI{0.00004}{\keV}$ \\  
$K=\frac{\SI{2.95}{\keV}}{127.264} = 0.023 \pm \SI{0.0003}{\keV}$ \\
Die Unsicherheiten wurden mithilfe der Gauß'schen Fehlerfortpflanzung berechnet. 
Den gesamten Kalibrierungsfaktor und seine Unsicherheit erhalten wir über Berechnung des gewichteten Mittelwerts und des zugehörigen Fehlers. Damit ergibt sich: \\ 
$K_{ges} = 0.024 \pm \SI{0.00004}{\keV}$ \\ 
Wenn man nun die pulseheight des Peaks von Cadmium und der x-ray gun mit diesem Faktor multipliziert, kann man erkennen, dass die resultierenden Werte sehr gut zu den aus der Literatur erwarteten Werte passen (Tabelle \ref{tab:Spektralanalyse}). Die Abweichung liegt etwa im 1$\sigma$-Intervall. 

\begin{table}
    \begin{tabular}{|l||c|c|c|}
        \hline 
        Quelle & Peakwert der pulseheight & Peak $\cdot K_{ges}$ & Literaturwert \\ 
        \hline 
        \hline 
        $^{55}$Fe & $246.822 \pm 0.438$ & & \\
        \hline 
        $ ^{55}$Fe (Argon-Escape) & $ 127.264 \pm 1.515$ & &  \\ 
        \hline 
        $^{109}$Cd & $343.020 \pm 4.177$ & $ 8.232 \pm $ \SI{0.101}{\keV} & \SI{8.34}{\keV} \cite{Anleitungsmappe}\\
        \hline 
        x-ray gun & $345.949 \pm 5.223$ & $8.30 \pm $ \SI{0.126}{\keV} & \SI{8.34}{\keV} \cite{Anleitungsmappe}\\
        \hline 
    \end{tabular}
    \label{tab:Spektralanalyse}
    \caption{Die aus dem Fit erhaltenen Werte des Peaks und die Multiplikation dieser mit dem Kalibrierungsfaktor, wobei die Unsicherheit auf diesen Wert mit Gauß'scher Fehlerfortpflanzung berechnet wurde.}
\end{table}

\begin{figure}[h]
    \centering
    \includesvg[width=0.8\textwidth]{Bilder/Spektren_Energie.svg}
    \caption{Mit dem Kalibrierungsfakotor multiplizierte Darstellung}
    \label{fig:Spektren_Energie}
    
\end{figure}
    
Das erhaltene Eisenspektrum zeigt zwei Peaks, was mit der Erwartung übereinstimmt. Der erste kleine Peak entspricht, wie oben schon erwähnt, dem Argon-Escape-Peak. Dieser entsteht, wenn ein Argon-Atom durch ein Röntgenphoton ionisiert wird. Meist wird dabei ein Elektron aus der K-Schale rausgeschlagen. Bei der anschließenden Rekombination wird ein Photon frei, was den Detektor auch undetektiert verlassen kann. Die Differenz der ursprünglichen Energie des Röntgenphotons und des K$_{\alpha}$-Photons sieht man dann als Peak im Spektrum. \\
Der zweite Peak entspircht dem Hauptpeak des $^{55}$Fe. \\ 
\\
Die Spektren der x-ray gun und des $^{109}$Cd sehen sehr ähnlich aus. Das lässt sich leicht erklären, denn die hauptsächlich emittierten Röntgenphotonen der x-ray gun und der $^{109}$Cd-Quelle haben den gleichen Wert von \SI{22.1}{\keV} \cite{Anleitungsmappe}. Nach der Absorption und Reemittierung aus der Kupferkathode haben die Photonen aus beiden Quellen die oben aufgeführte Energie von etwa \SI{8.34}{\keV}. Diesen Peak kann man in der Abbildung \ref{fig:Spektren_Energie} gut erkennen. Es scheint aber ebenfalls ein kleiner Peak um etwa \SI{5}{\keV} entstanden zu sein. Dies könnten ebenfalls Argon-Escape-Peaks sein. Die Energie der Ar K$_{\alpha}$-Linie beträgt \SI{2.96}{\keV} \cite{K_alpha-Linien}. Zieht man diese Energie von der Peakenergie ab erhält man \SI{5.38}{\keV}, was in etwa der Stelle des kleinen Peaks entspricht. \\


- Fe schmaler -> keine WW mit Kupferkathode? 
- Breite beeinflusst durch Stoßverbreiterung in Kupfer?
- Breite beeinflusst durch Bandstruktur von Kupfer? 

1. What are the differences between the sources?

5. How can you explain the peaks energy and their width based on physics you know?
6. What are the parameters that affect the peak width?
7. Would you use the GEM detector as a spectra analyzer? Why? 